\apendice{Plan de Proyecto \textit{Software}}


\section{Introducción}

La fase de planificación es una de las más críticas de cualquier proyecto
\textit{software}. En ella se estima el trabajo, el tiempo y el coste que puede
suponer su realización, para lo cual es necesario analizar con detalle cada
una de las partes que lo componen. Además de permitir estimar los recursos
necesarios, este análisis facilita y agiliza la toma de decisiones en las
etapas futuras del desarrollo.

Normalmente, la planificación de un proyecto se separa en dos partes:

\begin{itemize}
    \item A la parte que organiza el trabajo a lo largo del tiempo se le
    denomina \textbf{planificación temporal}.
    \item A la parte que evalúa si el proyecto es realista y factible
    se le denomina \textbf{estudio de viabilidad}.
\end{itemize}

En la planificación temporal se elabora un programa de tiempos en el que se
estima la duración de cada tarea del proyecto. Además, se fijan una \textbf{fecha
de inicio}, una \textbf{fecha de finalización estimada}, que se suele fijar
algo antes de la fecha límite para congelar el desarrollo de nuevas
funcionalidades y poner como objetivo probar el producto y detectar errores y,
por último, la \textbf{fecha límite}, que coincide con la fecha pactada con el
cliente para entregar el producto finalizado de acuerdo con los requisitos
indicados por este.

A la hora de estimar tiempos es importante no solo basarse en experiencias
previas con tareas similares, sino también tener en cuenta que ciertas
tareas pueden estar en una situación de bloqueo hasta que otras terminen.

Por otra parte, el estudio de viabilidad determina si completar el proyecto,
tanto a nivel legal, como dentro de los plazos establecidos con ciertos límites económicos es algo
realista, por lo tanto, se divide en dos apartados:

\begin{itemize}
    \item La parte que estima los costes que conlleva desarrollar el
    proyecto y sus posibles beneficios se conoce como \textbf{viabilidad
    económica}.
    \item La parte que analiza los apartados legales que pudiesen reportar
    conflictos, como por ejemplo las licencias de las diferentes herramientas
    que se van a usar, se conoce como \textbf{viabilidad legal}.


\end{itemize}

\section{Planificación temporal}
La entrega está prevista para el \fechaEntrega. El registro de trabajo
comienza el 26 de febrero de 2026, por lo que la estimación económica toma
como referencia los \num{204} días naturales comprendidos entre ambas
fechas, incluidas las dos. Las fechas y horas de los \textit{Sprints}
recogen el seguimiento realizado, mientras que la fecha de entrega fija el
límite previsto para completar el desarrollo y la documentación.

Para la planificación temporal se ha seguido una \textbf{metodología ágil}
adaptada a las características y necesidades del proyecto. El desarrollo
se ha dividido en iteraciones, denominadas \textit{Sprints}. Cada una de
ellas se estructuraba en torno a una reunión de equipo en la que se exponían
\textbf{los puntos tratados} durante dicha iteración, \textbf{los problemas}
que habían surgido a raíz del desarrollo de las tareas y, por último,
\textbf{los objetivos previstos} de cara al cierre de la siguiente iteración.

La duración de cada \textit{Sprint} ha sido de \textbf{2 semanas}, salvo en
algunas iteraciones concretas, en las que surgieron imprevistos.

Para organizar las tareas propuestas en cada \textit{Sprint} se ha usado un
\textbf{\textit{tablero Kanban}} en el que se iba marcando cuando se estaba
trabajando en una tarea o cuando se podía dar por finalizada.

En los siguientes apartados se detalla cada uno de los \textit{Sprints} por
separado.

\subsection{\textit{Sprint} 0: Planteamiento inicial e investigación}
\textbf{Duración:} del 26 de febrero al 12 de marzo de 2026.
Este \textit{Sprint} dio comienzo oficialmente al proyecto, aunque es cierto
que se realizó un estudio de viabilidad mínimo previo a este proyecto. En la
reunión inicial se explicó la idea más a fondo y se debatió qué dirección se
quería que tomase el proyecto desde el principio. Por tanto, este \textit{Sprint}
se centró en investigar y recopilar información útil para asentar una buena base
y en empezar a montar el principio del \textit{pipeline}, comenzando por el
\textbf{reconocimiento óptico de caracteres} (\acs{OCR}, del inglés \textit{Optical Character
Recognition})\acused{OCR}.

El desglose en tareas junto con los campos asociados a cada una se
recoge en la figura~\ref{fig:sprint-00}.

\imagen{sprint-00}{Tablero del \textit{Sprint} 0.}

Se emplearon 13 horas y media de las 17 estimadas, y se completaron todas las
tareas previstas.

\subsection{\textit{Sprint} 1: Arquitectura base, CI y módulo OCR}
\textbf{Duración:} del 12 de marzo al 26 de marzo de 2026.
Con una idea más asentada acerca de la dirección del proyecto, este \textit{Sprint}
se centró en configurar el repositorio y su \textbf{integración continua} (CI, del
inglés \textit{Continuous Integration})\acused{CI}. Además, se creó una base de
arquitectura sobre la que empezar a desarrollar las funcionalidades.
Por otro lado, se investigaron y probaron diferentes opciones de \ac{OCR}, de
\textbf{generación aumentada por recuperación}
(\acs{RAG}, del inglés \textit{Retrieval-Augmented Generation})\acused{RAG}
y de \textbf{grandes modelos de lenguaje}
(\acs{LLM}, del inglés \textit{Large Language Model})\acused{LLM}.

El desglose en tareas junto con los campos asociados a cada una se
recoge en la figura~\ref{fig:sprint-01}.

\imagen{sprint-01}{Tablero del \textit{Sprint} 1.}

Se emplearon 24 horas de las 24 estimadas, y se completaron todas las
tareas previstas.

\subsection{\textit{Sprint} 2: Migración a microservicios, RAG y LLM}
\textbf{Duración:} del 26 de marzo al 16 de abril de 2026.
Este \textit{Sprint} se centró en pulir la estructura base del proyecto
y migrar la arquitectura monolítica a una arquitectura de
microservicios, en la que se planteó mover cada elemento lo
suficientemente grande, como por ejemplo, el \ac{LLM} o el \ac{OCR},
a su propio módulo en el que poder modificarse o instalar dependencias
sin afectar al resto del sistema~\cite{fowler_microservices}.
Adicionalmente, se hizo uso de esta nueva arquitectura para montar la
primera versión del \textit{pipeline}, conectando el \ac{RAG} con el \ac{LLM}.

El desglose en tareas junto con los campos asociados a cada una se
recoge en la figura~\ref{fig:sprint-02}.

\imagen{sprint-02}{Tablero del \textit{Sprint} 2.}

Se emplearon 28 horas frente a las 20 estimadas, y se completaron todas las
tareas previstas.

\subsection{\textit{Sprint} 3: Estabilización y optimización del motor de IA}
\textbf{Duración:} del 16 de abril al 23 de abril de 2026.

Este \textit{Sprint} se centró en mejorar la estabilidad y optimizar el motor de
\ac{IA} de la aplicación, formado por el \ac{RAG} y el \ac{LLM}. Para ello, se
establecieron una serie de requisitos y se siguió un enfoque de desarrollo guiado
por pruebas (\acs{TDD}, del inglés \textit{Test-Driven Development})\acused{TDD}~\cite{agile_alliance_tdd},
con el objetivo de validar el comportamiento esperado de dichos módulos.

Por otro lado, se hizo un análisis del estado del arte de \textbf{modelos de
síntesis de voz} (\acs{TTS}, del inglés \textit{Text-To-Speech})\acused{TTS} y un
estudio de viabilidad mínimo, sobre todo de cara a las limitaciones a nivel de
\textit{hardware} disponible.

El desglose en tareas junto con los campos asociados a cada una se
recoge en la figura~\ref{fig:sprint-03}.

\imagen{sprint-03}{Tablero del \textit{Sprint} 3.}

Se emplearon 15 horas frente a las 12 estimadas, y se completaron todas las
tareas previstas.

\subsection{\textit{Sprint} 4: Creación de un \textit{frontend} básico funcional}
\textbf{Duración:} del 23 de abril al 12 de mayo de 2026.

Este \textit{Sprint} se centró en refactorizar el \textit{backend}, aplicando
patrones de diseño y buenas prácticas. Por ejemplo, se emplearon patrones como \textbf{\textit{Strategy}} y \textbf{\textit{Factory}}~\cite{patterns}. El
primero sirve para encapsular los distintos motores de \ac{OCR} tras una interfaz
común e intercambiarlos sin afectar al resto del sistema, y el segundo para
instanciar el motor adecuado según la configuración.

Por otro lado, se implementaron en el flujo de \ac{CI} herramientas como
\textbf{\textit{Ruff}}~\cite{ruff} para mejorar la calidad del código.

Sobre esa base, se conectó una interfaz sencilla en \acs{HTML} (del inglés
\textit{HyperText Markup Language})\acused{HTML} con la que disponer de una demo y probar el \textit{pipeline} de \textbf{extremo a
extremo} (\acs{E2E}, del inglés \textit{End-to-End})\acused{E2E}.

La primera interfaz web de demostración se muestra en la figura~\ref{fig:sprint-04-demo-frontend}.

\imagen{sprint-04-demo-frontend}{Interfaz web de demostración desarrollada para probar el \textit{pipeline} de extremo a extremo.}

En la reunión se debatió acerca del estudio de viabilidad del módulo de \ac{TTS}
y, finalmente, se decidió dejarlo fuera debido a las limitaciones de
\textit{hardware} y a la complejidad y la pérdida de calidad en
la experiencia de usuario que supondría un sistema de carga y descarga de modelos
para gestionar los recursos.

El desglose en tareas junto con los campos asociados a cada una se
recoge en la figura~\ref{fig:sprint-04}.

\imagen{sprint-04}{Tablero del \textit{Sprint} 4.}

Se emplearon 20 horas y 40 minutos frente a las 20 estimadas, y se
completaron todas las tareas previstas.

\subsection{\textit{Sprint} 5: Refinamiento del \textit{backend} y diseño del \textit{frontend}}
\textbf{Duración:} del 12 de mayo al 26 de mayo de 2026.

Este \textit{Sprint} se centró principalmente en desarrollar la primera versión de
una \textbf{aplicación web progresiva} (\acs{PWA}, del inglés
\textit{Progressive Web App})\acused{PWA} y sentar las bases para mejorarla en
futuras iteraciones~\cite{mdn_pwa}.

Adicionalmente, como resultado de haber realizado la demo, se observó que, pese a
que el motor \textit{\textbf{PaddleOCR}}~\cite{paddleocr_tech_report} que se estaba
utilizando era muy fiable en cuanto a los resultados, era considerablemente lento. Además,
su variante adaptada para la \textbf{\textit{Unidad de Procesamiento Gráfico}}
(\acs{GPU}, del inglés \textit{Graphics Processing Unit})\acused{GPU}, con la que se reducía el
problema de rendimiento, es muy dependiente de las versiones de \acs{CUDA} (del inglés
\textit{Compute Unified Device Architecture})\acused{CUDA}, \textit{PaddlePaddle} y de los
\textit{drivers} de la \ac{GPU}~\cite{paddlepaddle_install}. Por motivos de compatibilidad,
se decidió añadir el motor \ac{OCR} \textbf{\textit{Tesseract}}~\cite{tesseract_ocr} como
opción por defecto y haciendo uso de los patrones de diseño aplicados que permiten tener
varias opciones configurables~\cite{patterns}, tras haber sido el mejor motor en promedio
en las medidas tomadas en los \textit{benchmarks} realizados.


Por último, se implementó \textit{SonarQube}~\cite{sonarqube_ci} como un paso más de
\ac{CI} con el objetivo de
reducir \textit{code smells}~\cite{sonarqube_metrics} y malas prácticas en el código.
Utilizando esta nueva herramienta, en este mismo \textit{Sprint} se aprovechó para
corregir la \textbf{deuda técnica}~\cite{sonarqube_metrics} detectada por
\textit{SonarQube}.

El desglose en tareas junto con los campos asociados a cada una se
recoge en la figura~\ref{fig:sprint-05}.
\imagen{sprint-05}{Tablero del \textit{Sprint} 5.}

Se emplearon 81 horas y 7 minutos frente a las 58 estimadas, y se
completaron todas las tareas previstas.

\subsection{\textit{Sprint} 6: Finalización del MVP}
\textbf{Duración:} del 26 de mayo al 13 de junio de 2026.

Este \textit{Sprint} se centró en montar \textbf{\textit{una base de datos relacional en
PostgreSQL}}~\cite{postgresql_about}, sobre la que almacenar manuales, \textit{assets},
conversaciones, etc.

La parte más importante de este paso se encontraba en crear un \textbf{sistema de usuarios},
haciendo especial énfasis en su \textbf{seguridad} teniendo en cuenta apartados
como \textbf{la gestión de sesiones}~\cite{owasp_session_management}, o la \textbf{protección
contra falsificación de peticiones entre sitios web} (\acs{CSRF}, del inglés \textit{Cross-Site
Request Forgery})\acused{CSRF}~\cite{owasp_csrf} entre otros.
Estas medidas se detallan en el apartado
\guillemotleft{}\nameref{sec:diseno-seguridad}\guillemotright{} del anexo de diseño.

Al mismo tiempo, se puso como objetivo mejorar el \textit{frontend} y la experiencia de
usuario mediante iteraciones hasta llegar a tener un \textbf{producto mínimo viable}
(\acs{MVP}, del inglés \textit{Minimum Viable Product})\acused{MVP}~\cite{agile_alliance_mvp}, incluyendo
apartados como el rendimiento, la flexibilidad y la robustez, tal como se puede ver en el
desglose de tareas al ampliar la cantidad de \textit{assets} y formatos soportados o al reducir
el tamaño y tiempo de construcción de las imágenes \textit{Docker}.

El desglose en tareas junto con los campos asociados a cada una se
recoge en la figura~\ref{fig:sprint-06}.

\imagen{sprint-06}{Tablero del \textit{Sprint} 6.}

Se emplearon 151 horas y 40 minutos frente a las 77 estimadas, y se
completaron todas las tareas previstas.

\subsection{\textit{Sprint} 7: Preparación del despliegue}
\textbf{Duración:} del 15 de junio al 30 de junio de 2026.

Este \textit{Sprint} se centró sobre todo en dejar todo preparado para el despliegue
de la aplicación. Para ello, se llevó a cabo una fase de
\textbf{aseguramiento de la calidad} (\acs{QA}, del inglés
\textit{Quality Assurance})\acused{QA}~\cite{iso_quality_assurance}.

Durante dicha fase, se detectaron principalmente tres problemas:

\begin{itemize}
    \item El procesamiento de tareas pesadas estaba bloqueando el
    \textbf{bucle de eventos principal de la aplicación}.
    \item La calidad de los resultados del \ac{OCR} era, en ocasiones,
    \textbf{insuficiente} para enriquecer las respuestas del \ac{LLM}
    \item La aplicación requería de un \textit{hardware} muy específico.
\end{itemize}

Una vez identificados, se crearon tareas para investigar posibles soluciones,
implementarlas y comprobar que corregían los problemas detectados.

La ejecución tareas pesadas se movieron a una \textbf{cola de Celery}, utilizando Redis como \textit{broker} de mensajes~\cite{celery_introduction}.

Para mejorar los resultados del \ac{OCR}, se creó un \textit{pipeline} que
aplica un \textbf{preprocesado de imagen} específico para cada motor.

Por último, se añadió un arranque adaptativo que detecta la disponibilidad de
una GPU NVIDIA y la VRAM libre para recomendar perfiles compatibles del \ac{LLM}
y del \ac{OCR}.

El desglose en tareas junto con los campos asociados a cada una se
recoge en la figura~\ref{fig:sprint-07}.

\imagen{sprint-07}{Tablero del \textit{Sprint} 7.}

Se emplearon 82 horas frente a las 82 estimadas, y se completaron todas
las tareas previstas.

\subsection{\textit{Sprint} 8: Despliegue de la aplicación}
\textbf{Duración:} del 1 de julio al 20 de julio de 2026.

Este \textit{Sprint} se centró en crear la infraestructura necesaria para
realizar el primer despliegue de \textit{Manualito}, tanto del sitio web de
presentación (desplegado en \href{https://manualito.dev/}{\textit{manualito.dev}})
como de la aplicación a través de \textit{Cloudflare
Tunnel}~\cite{cloudflare_tunnel}
(desplegado a demanda en \href{https://app.manualito.dev/}{\textit{app.manualito.dev}}).

Como consecuencia de preparar el proyecto para su publicación, se estableció un
flujo de trabajo basado en \textbf{ramas por entorno}~\cite{gitlab_branch_per_environment}
para organizar el desarrollo, las pruebas y el despliegue de los cambios.

Además, los \textit{runners} de GitLab alojados en la infraestructura de HP SCDS
comenzaron a fallar, lo que impedía mantener el flujo de \ac{CI} e implementar un
procedimiento fiable de \textbf{despliegue continuo} (\acs{CD}, del inglés
\textit{Continuous Deployment})\acused{CD}. Por este motivo, el repositorio
principal se migró a GitHub~\cite{manualito_github}. Por su parte, GitLab se
mantuvo para conservar la trazabilidad del proyecto, por lo que las
\textit{issues} se replicaron en ambas plataformas~\cite{manualito_gitlab}.

El desglose en tareas junto con los campos asociados a cada una se
recoge en la figura~\ref{fig:sprint-08}.

\imagen{sprint-08}{Tablero del \textit{Sprint} 8.}

Se emplearon 68 horas y 35 minutos frente a las 75 estimadas, y se
completaron todas las tareas previstas.

\subsection{\textit{Sprint} 9: Despliegue de la documentación}
\textbf{Duración:} del 20 de julio al 6 de agosto de 2026.

Este \textit{Sprint} se centró en crear un sitio web en el que publicar la
documentación de \textit{Manualito}. Para ello, se utilizó \textit{Starlight},
un tema de documentación basado en \textit{Astro}~\cite{starlight_documentation}, para exponer las guías de uso y la documentación técnica del proyecto.

El sitio se desplegó mediante
\textit{Cloudflare Workers Static Assets}~\cite{cloudflare_workers_static_assets}
y se publicó en \href{https://docs.manualito.dev/}{\textit{docs.manualito.dev}}.

Por otro lado, se internacionalizó la interfaz de \textit{Manualito} para que
pudiera utilizarse \textbf{tanto en español como en inglés}.

El desglose en tareas junto con los campos asociados a cada una se
recoge en la figura~\ref{fig:sprint-09}.

\imagen{sprint-09}{Tablero del \textit{Sprint} 9.}

Se emplearon 50 horas y 30 minutos frente a las 53 estimadas, y se
completaron todas las tareas previstas.

\subsection{\textit{Sprint} 10: Mejora del \acs{RAG}}
\textbf{Duración:} del 6 de agosto al 23 de agosto de 2026.

Este \textit{Sprint} se centró en mejorar el \ac{RAG} de \textit{Manualito}.
Para ello, se combinaron las \textbf{búsquedas semánticas y textuales}, y se
incorporó el nombre del juego a la consulta para interpretar preguntas que
no lo mencionaban explícitamente.

Además, se aplicó el filtro de \textbf{manuales autorizados} antes de
recuperar los fragmentos, para que los documentos que el usuario no podía
consultar no ocupasen posiciones entre los resultados.

Por último, se añadió una tarea periódica para mantener el índice de
búsqueda sincronizado con PostgreSQL, eliminando las entradas de manuales
borrados y repitiendo las indexaciones incompletas.

El desglose en tareas junto con los campos asociados a cada una se
recoge en la figura~\ref{fig:sprint-10}.

\imagen{sprint-10}{Tablero del \textit{Sprint} 10.}

Se emplearon 47 horas frente a las 115 estimadas, y se completaron todas las tareas previstas.

\subsection{\textit{Sprint} 11: Redacción de la documentación}
\textbf{Duración:} del 23 de agosto al 17 de septiembre de 2026.

Este \textit{Sprint} se centró en recopilar y completar la documentación
redactada durante el desarrollo para preparar la \textbf{memoria y los anexos
del trabajo de fin de grado}. También se publicaron los estudios de rendimiento
del proyecto y se corrigieron pequeños errores e incoherencias detectados al
revisar la interfaz y el envío de correos.

El desglose en tareas junto con los campos asociados a cada una se
recoge en la figura~\ref{fig:sprint-11}.

\imagen{sprint-11}{Tablero del \textit{Sprint} 11.}

El tiempo dedicado a la memoria y los anexos se repartió entre numerosas sesiones
breves, por lo que no se conserva una medición suficientemente precisa de las
horas empleadas. Para evitar establecer una cifra reconstruida sin evidencia,
no se incluye una estimación del esfuerzo de este \textit{Sprint}.

\section{Estudio de viabilidad}

Una vez planificado el trabajo en el tiempo, queda comprobar que el proyecto
se puede realizar. Para ello se analizan los dos apartados que mencionados
anteriormente: la \textbf{viabilidad económica}, que estima los costes de
desarrollo y los contrasta con los beneficios esperables, y la \textbf{viabilidad
legal}, que revisa las licencias de las herramientas empleadas y el cumplimiento
normativo aplicable.

\subsection{Viabilidad económica}
El presupuesto estima cuánto costaría desarrollar \textit{Manualito} en un
entorno empresarial, aunque el proyecto se haya realizado en un contexto
académico. Para valorar el trabajo se mantiene la hipótesis de un
desarrollador a jornada completa durante el periodo previsto, del 26 de
febrero al \fechaEntrega.

El personal, las amortizaciones y la conexión a internet se prorratean mediante el factor
\(\num{204}/\num{365}\), correspondiente a los días del periodo y del año
2026. Las cuotas mensuales se convierten primero a su equivalente anual.
Así se incluyen los meses parciales sin redondear la duración a meses
completos. Los importes se redondean a céntimos antes de sumar las filas de
cada tabla.

\subsubsection{Costes}
La estimación reúne el coste de personal, la parte del equipo y las licencias
imputada al proyecto, y los gastos de electricidad, conexión y entrega.

\noindent\textbf{Costes de personal}\par
Se considera un único desarrollador de perfil \textit{junior} a jornada
completa. Como referencia mínima se toma el \textbf{Salario Mínimo
Interprofesional} de 2026, fijado en \num{17094.00}\,€ anuales~\cite{smi2026}.
Esta cantidad sirve como base del presupuesto, sin equipararla al salario de
mercado de ese perfil.

Sobre esa base se estima una cotización empresarial del
\qty{30.65}{\percent}, suponiendo una contratación indefinida. El porcentaje
reúne las contingencias comunes (\qty{23.60}{\percent}), el desempleo
(\qty{5.50}{\percent}), el Mecanismo de Equidad Intergeneracional
(\qty{0.75}{\percent}), el Fondo de Garantía Salarial
(\qty{0.20}{\percent}) y la formación profesional
(\qty{0.60}{\percent})~\cite{orden_cotizacion_2026}.

Se trata de una aproximación presupuestaria. Una contratación real exigiría
comprobar la base mínima del grupo profesional y añadir la cotización por
accidentes de trabajo y enfermedades profesionales, cuya tarifa depende de
la actividad. Tampoco se contemplan horas extraordinarias ni complementos
salariales.

La tabla~\ref{tabla:costes-personal} recoge la referencia anual y el importe
imputado a los \num{204} días del proyecto.
\tablaSmallSinColores
  {Costes de personal.}
  {l S[table-format=5.2, detect-weight]}
  {costes-personal}
  {\textbf{Concepto} & {\textbf{Coste (€)}} \\}
  {
    Salario bruto anual                       & \num{17094.00} \\
    Cotización empresarial anual              & \num{5239.31} \\
    Coste empresarial anual estimado           & \num{22333.31} \\
    \midrule
    Salario bruto del periodo                 & \num{9553.91} \\
    Cotización empresarial del periodo        & \num{2928.27} \\
    \midrule
    \bfseries Total (204 días)                & \bfseries \num{12482.18} \\
  }

\noindent\textbf{Costes de componentes \textit{hardware}}\par

 El siguiente paso es revisar todos los gastos derivados de componentes \textit{hardware}
 que se necesitan para desarrollar y probar un producto \textit{software} de este tipo. En
 este caso es especialmente relevante, ya que a la hora de tratar con modelos \ac{IA}
 como los \ac{LLM} se necesitan ciertas capacidades de computación que idealmente se
 obtienen de una \ac{GPU} que cuente con la suficiente potencia.

A continuación se desglosan los componentes y costes asociados que podría tener un equipo
con la potencia necesaria~\ref{tabla:costes-hardware}.

\tablaSmallSinColores
  {Costes de componentes \textit{hardware}.}
  {l r r}
  {costes-hardware}
  {\textbf{Componente} & \textbf{Precio (€)} & \textbf{Amortizado (€)} \\}
  {
    Procesador AMD Ryzen 5 7600             & \num{174.90}  & \num{24.44}  \\
    \ac{GPU} NVIDIA RTX 4070 Super          & \num{639.90}  & \num{89.41}  \\
    Memoria RAM 32\,GB DDR5                 & \num{184.98}  & \num{25.85}  \\
    Placa base AM5 (B850)                   & \num{129.90}  & \num{18.15}  \\
    SSD NVMe PCIe 4.0 1\,TB                 & \num{154.95}  & \num{21.65}  \\
    Fuente 750\,W 80+ Gold                  & \num{112.99}  & \num{15.79}  \\
    Disipador por aire AM5                  & \num{21.95}   & \num{3.07}   \\
    Caja ATX                                & \num{79.90}   & \num{11.16}   \\
    \midrule
    \bfseries Total                         & \bfseries\num{1499.47} & \bfseries\num{209.52} \\
  }

Los precios se han consultado en PcComponentes a fecha de 20 de junio de 2026 ~\cite{pccomp_ryzen_7600,
pccomp_rtx_4070_super,pccomp_ram_kingston_32gb,pccomp_msi_b850,
pccomp_kingston_nv3,pccomp_msi_a750gl,pccomp_hyper_212,pccomp_corsair_3000d}.

Para estimar la parte del equipo consumida durante el proyecto se adopta
una amortización lineal con una vida útil de cuatro años (\num{48} meses),
valor residual nulo e imputación íntegra durante el periodo considerado.
Esta vida útil equivale al coeficiente anual del \qty{25}{\percent}
recogido para los equipos informáticos en la tabla del impuesto sobre
sociedades~\cite{ley_impuesto_sociedades}, que se utiliza como referencia.
El coste de cada componente se calcula como:
\[
  \text{Coste amortizado} = \text{Precio} \times
  \frac{\num{204}}{\num{365} \times \num{4}}
\]

\noindent\textbf{Costes de licencias \textit{software}}\par

Otro punto importante es revisar los gastos en licencias de programas, si bien la
gran mayoría de \textit{software} utilizado es de código abierto (\acs{OSS}, del
inglés \textit{Open Source Software})\acused{OSS}, el desarrollo y las pruebas
se han realizado sobre \textit{Windows}, lo que requiere una licencia de
\textit{Windows 11 Pro}~\cite{ms_windows11_pro}.

Los precios de las licencias de \textit{software} se desglosan en la siguiente
tabla~\ref{tabla:costes-software}.

\tablaSmallSinColores
  {Costes de las licencias \textit{software}.}
  {l S[table-format=3.2] S[table-format=2.2]}
  {costes-software}
  {\textbf{\textit{Software}} & {\textbf{Precio (€)}} & {\textbf{Amortizado (€)}} \\}
  {
    Windows 11 Pro                         & \num{259.00} & \num{48.25} \\
    \midrule
    \bfseries Total                        & \bfseries \num{259.00} & \bfseries \num{48.25} \\
  }

Para la licencia se adopta una vida útil de tres años (\num{36} meses),
con amortización lineal y valor residual nulo. Es una hipótesis del
presupuesto, por lo que el importe imputado se obtiene mediante
\(\num{259.00} \times \num{204}/(\num{365} \times \num{3})\), que resulta
en \num{48.25}\,€.

\noindent\textbf{Otros costes}\par
El consumo eléctrico se aproxima a partir del esfuerzo registrado. Los
\textit{Sprints} 0 a 9 suman \num{535} horas y \num{2} minutos entre el
26 de febrero y el 6 de agosto, un periodo de \num{162} días naturales.
Para estimar el consumo hasta la entrega se mantiene ese ritmo medio
durante los \num{204} días previstos:
\[
  \text{Horas estimadas} =
  \left(\num{535} + \frac{\num{2}}{\num{60}}\right)
  \times \frac{\num{204}}{\num{162}}
  \simeq \num{673.75}
\]

Suponiendo que el equipo se utiliza durante esas horas, con una potencia
media de \qty{250}{\watt} y un precio de \num{0.15}\,€/kWh~\cite{ocu_precio_luz},
el coste estimado es de \num{25.27}\,€. De ese importe, \num{20.06}\,€
corresponden al esfuerzo ya registrado y \num{5.21}\,€ a la prolongación del
mismo ritmo hasta la entrega. Esta aproximación no equivale a una medición
del consumo ni añade horas al registro de los \textit{Sprints}.

Para la conexión a internet se conserva una cuota de
\num{30}\,€ mensuales~\cite{jazztel_fibra} y se imputa al proyecto un
\qty{30}{\percent} del uso, al compartirla con actividades personales.
Aplicando el mismo periodo que al resto del presupuesto, el cálculo es
\(\num{30} \times \num{12} \times \num{204}/\num{365}
\times \num{0.30} = \num{60.36}\)\,€.

La impresión y encuadernación se estiman en \num{50}\,€ como gasto puntual.
Para el registro del dominio \textit{manualito.dev} se reserva una anualidad
de \num{12}\,€.


Todo esto se desglosa en la tabla~\ref{tabla:otros-costes}.

\tablaSmallSinColores
  {Otros costes.}
  {l S[table-format=3.2]}
  {otros-costes}
  {\textbf{Concepto} & {\textbf{Coste (€)}} \\}
  {
    Consumo eléctrico                        & \num{25.27} \\
    Conexión a internet                      & \num{60.36} \\
    Impresión y encuadernación de la memoria & \num{50.00} \\
    Dominio web                              & \num{12.00} \\
    \midrule
    \bfseries Total                          & \bfseries \num{147.63} \\
  }

\noindent\textbf{Costes operativos}\par
Además de los costes de desarrollo, mantener la aplicación disponible de
forma continuada tendría asociados costes operativos mensuales.
Estos costes dependerían principalmente del uso real: número de usuarios
activos, mensajes procesados, modelo \ac{LLM}, almacenamiento de manuales y \textit{assets} e infraestructura de despliegue.

Durante el proyecto se ha utilizado el plan gratuito de Resend, sin coste
alguno. El envío está sujeto a los límites diarios y mensuales de ese
plan~\cite{resend_precios}.

Por este motivo, no se fija una cifra cerrada para dichos costes en esta fase del
proyecto. Al depender directamente del uso real de la aplicación, deberían
estimarse mediante pruebas de carga y métricas de uso antes de ofrecer el
servicio de forma continuada.

Hay que tener en cuenta que estos costes condicionarían tanto los límites de uso
de cada plan como el precio final de las suscripciones.

\noindent\textbf{Coste efectivo del proyecto}\par
Una vez estimados los costes de personal, \textit{hardware}, \textit{software} y
otros gastos asociados, se obtiene el coste imputado al desarrollo del proyecto.
En este cálculo no se incluye el coste completo de adquisición del equipo, sino
únicamente la parte amortizada correspondiente a los \num{204} días previstos,
ya que dicho equipo podría seguir utilizándose una vez finalizado el proyecto.

El total de coste imputado al proyecto se desglosa en la
tabla~\ref{tabla:coste-imputado-proyecto}.

\tablaSmallSinColores
  {Coste imputado al proyecto.}
  {l r}
  {coste-imputado-proyecto}
  {\textbf{Concepto} & \textbf{Coste (€)} \\}
  {
    Costes de personal        & \num{12482.18} \\
    Costes de \textit{hardware}        & \num{209.52}  \\
    Costes de \textit{software}        & \num{48.25}    \\
    Otros costes              & \num{147.63}  \\
    \midrule
    \bfseries Total           & \bfseries\num{12887.58} \\
  }

Con este dato, ya podemos estimar los beneficios que podemos obtener con el producto.

\subsubsection{Beneficios}

Esta aplicación \textbf{no está pensada para ser comercializada}, principalmente porque su
código fuente ha sido publicado y, adicionalmente, se han detallado los mecanismos de seguridad
que se usan en ella. Por este motivo, sería imprudente, no solo comercializar la aplicación
a corto plazo, sino incluso desplegarla fuera del contexto académico con el objetivo de
utilizar datos reales de usuarios.

Dicho esto y una vez se han comprendido los riesgos, sí que se podría llegar a comercializar
la aplicación en una segunda fase, tras haber cambiado y mejorado los sistemas actuales de
seguridad para corregir vulnerabilidades y tras haber pasado las auditorías de seguridad
necesarias.

En este caso, se distribuiría bajo un modelo de \textbf{\textit{software} como servicio}
(\acs{SaaS}, del inglés \textit{Software as a Service})\acused{SaaS}~\cite{nist_cloud_computing} en el que se ofrecerían planes tanto para acceder a mejores modelos \ac{LLM}, como para aumentar
los límites de mensajes diarios por cuenta. En cualquier caso, los datos y \textit{assets} se
guardarían en la nube.

Si se llegase a comercializar se plantearían tres tipos de suscripciones mensuales detallados
en la tabla~\ref{tabla:suscripciones}.

\tablaSmallSinColores
  {Planes de suscripción.}
  {l c c c}
  {suscripciones}
  { & \textbf{Aprendiz} & \textbf{Estratega} & \textbf{Gran Maestro} \\}
  {
    Precio mensual       & Gratis        & \num{4.99}\,€ & \num{9.99}\,€  \\
    Juegos patrocinados  & Sí            & No            & No             \\
    Modelo \ac{LLM}      & granite3.3:2b & gemma4:e4b    & gemma4:e4b     \\
    Mensajes diarios     & \num{20}      & \num{200}     & Ilimitados     \\
    Manuales             & \num{5}       & \num{50}      & Ilimitados     \\
    Almacenamiento       & \num{100}\,MB & \num{2}\,GB   & \num{20}\,GB   \\
    Soporte              & Comunidad     & Correo        & Prioritario    \\
  }




\subsection{Viabilidad legal}
\label{sec:viabilidad-legal}
\begingroup
\DefTblrTemplate{conthead-text}{legal}{(continuación)}
\DefTblrTemplate{contfoot-text}{legal}{Continúa en la página siguiente}
\SetTblrTemplate{conthead-text,contfoot-text}{legal}

\textit{Manualito} combina código propio con bibliotecas, modelos y recursos
visuales de terceros, cuyo uso y distribución están sujetos a las condiciones
de sus respectivas licencias.

\subsubsection{\textit{Software}}

El código de \textit{Manualito} se publica bajo la \textbf{licencia \acs{MIT}},
recogida en el archivo \texttt{LICENSE} del repositorio
\cite{manualito_licencia}. Esta licencia permite utilizar, estudiar, modificar
y distribuir el programa, incluso con fines comerciales, siempre que se
conserven el aviso de autoría y el texto de la licencia en las copias o partes
sustanciales del código. De este modo, otras personas pueden adaptar la
aplicación a sus necesidades sin tener que solicitar un permiso individual
para cada uso~\cite{licencia_mit}.

La elección de \acs{MIT} se aplica al código propio y no sustituye las licencias de
las dependencias. Para conocer qué condiciones acompañan a la aplicación se
han revisado las versiones declaradas en el proyecto y sus fuentes oficiales.
La tabla~\ref{tabla:licencias-componentes} recoge los componentes principales,
con su función y la licencia de cada componente. En las
bibliotecas que incorporan otros componentes, la licencia principal se
completa con los avisos de esos materiales.

\begin{longtblr}[
  caption = {Licencias de los componentes principales.},
  label = {tabla:licencias-componentes},
]{
  width = \linewidth,
  colspec = {X[2.8,l] X[1.5,l] X[3.0,l] X[2.7,l]},
  rowhead = 1,
  cells = {font=\small},
  row{1} = {font=\small\bfseries},
  rowsep = 3pt,
}
\toprule
Componente & Versión & Uso & Licencia \\
\midrule
FastAPI & 0.141.1 & Atención de peticiones. & \acs{MIT}~\cite{licencia_fastapi} \\
React & 19.2.8 & Interfaz del navegador. & \acs{MIT}~\cite{licencia_react} \\
SQLAlchemy & 2.0.51 & Acceso a los datos. & \acs{MIT}~\cite{licencia_sqlalchemy} \\
Psycopg & 3.3.4 & Comunicación con PostgreSQL. & \acs{LGPL} 3.0~\cite{licencia_psycopg} \\
PostgreSQL & 18.4 & Persistencia de datos. & PostgreSQL~\cite{licencia_postgresql} \\
Redis & 8.10.0 & Colas y coordinación de tareas. & \acs{AGPL} 3.0, \acs{RSAL} 2 o \acs{SSPL} 1, a elección~\cite{licencia_redis} \\
Celery & 5.6.3 & Tareas en segundo plano. & \acs{BSD}, tres cláusulas~\cite{licencia_celery} \\
Chroma & 1.5.9 & Índice de vectores. & Apache 2.0~\cite{licencia_chroma} \\
Sentence Transformers & 5.7.0 & Representación vectorial del texto. & Apache 2.0~\cite{licencia_sentence_transformers} \\
Ollama & 0.32.6 & Ejecución de los modelos de lenguaje. & \acs{MIT}~\cite{licencia_ollama} \\
PaddleOCR & 3.7.0 & Reconocimiento de texto. & Apache 2.0~\cite{licencia_paddleocr} \\
PaddlePaddle & 3.3.1 & Modelos de PaddleOCR. & Apache 2.0 para la biblioteca~\cite{licencia_paddlepaddle} \\
Tesseract & No fijada & Reconocimiento alternativo de texto. & Apache 2.0~\cite{licencia_tesseract} \\
pypdfium2 & 5.12.1 & Lectura de documentos \ac{PDF}. & Apache 2.0 o \acs{BSD}, tres cláusulas, para su código~\cite{licencia_pypdfium2} \\
Pillow & 12.3.0 & Lectura de imágenes. & MIT-CMU~\cite{licencia_pillow} \\
OpenCV, paquete Python & 4.10.0.84 & Preparación de imágenes. & Apache 2.0 para OpenCV y \acs{MIT} para su empaquetado~\cite{licencia_opencv} \\
Caddy & 2.11.4 & Servidor de la interfaz y \textit{proxy} inverso. & Apache 2.0~\cite{licencia_caddy} \\
Mailpit & 1.30.6 & Servidor de correo local. & \acs{MIT}~\cite{licencia_mailpit} \\
Flower & 2.0.1 & Supervisión de tareas. & \acs{BSD}, tres cláusulas~\cite{licencia_flower} \\
\bottomrule
\end{longtblr}

La versión de Tesseract depende del paquete instalado al construir la imagen,
por lo que debe concretarse junto con la versión final que se entregue.
En la fila de OpenCV se indica la versión del paquete
\texttt{opencv-contrib-python} que utiliza la aplicación.

Las licencias \acs{MIT} y Apache, junto con las variantes
\textbf{\acs{BSD}} (del inglés \textit{Berkeley Software Distribution})\acused{BSD},
permiten reutilizar y distribuir las bibliotecas conservando sus avisos.
Apache 2.0 exige también señalar los
archivos modificados y mantener las atribuciones de un archivo
\texttt{NOTICE} cuando forme parte de la distribución original
\cite{licencia_mit,licencia_apache,licencia_bsd}. Por tanto, publicar el código
de \textit{Manualito} bajo \acs{MIT} es compatible con mantener esos componentes
bajo sus propias condiciones.

Algunas bibliotecas añaden requisitos cuando se entregan dentro de la
aplicación. Por ejemplo, Psycopg utiliza la
\textbf{licencia pública general reducida}
(\acs{LGPL}, del inglés \textit{GNU Lesser General Public License})\acused{LGPL},
que permite combinar la biblioteca con código de otra
licencia, pero exige conservar sus avisos y cumplir las condiciones de acceso
al código y sustitución de la biblioteca que correspondan al modo de
distribución~\cite{licencia_psycopg}. Esta distinción importa si se entrega una
imagen Docker con las dependencias ya instaladas, porque no contiene solo el
código propio. OpenCV y PDFium también incorporan materiales con avisos
adicionales, que deben acompañar a los componentes que se distribuyan
\cite{licencia_opencv,licencia_pypdfium2}.

Redis 8 permite escoger entre \textit{GNU Affero General Public License}
(\acs{AGPL})\acused{AGPL}, \textit{Redis Source Available License}
(\acs{RSAL})\acused{RSAL} y \textit{Server Side Public License}
(\acs{SSPL})\acused{SSPL}. En \textit{Manualito} se utiliza
como un servicio independiente, de modo que su presencia no obliga por sí
sola a cambiar la licencia del código propio. No obstante, la opción elegida
debe respetarse al distribuir Redis o modificarlo, y el uso de una versión
modificada a través de la red puede añadir obligaciones de acceso a su código
\cite{licencia_redis}.

\noindent\textbf{Modelos}\par

Los modelos que ejecuta la aplicación tienen condiciones distintas de las
bibliotecas que los cargan. Por ello, la licencia \acs{MIT} de Ollama no determina
la de los modelos de lenguaje, del mismo modo que la licencia de Sentence
Transformers no sustituye la del modelo utilizado para representar el texto.
La tabla~\ref{tabla:licencias-modelos} recoge los modelos configurados para
estas funciones.

\begin{longtblr}[
  caption = {Licencias de los modelos configurados en la aplicación.},
  label = {tabla:licencias-modelos},
]{
  width = \linewidth,
  colspec = {X[3.2,l] X[4.0,l] X[2.8,l]},
  rowhead = 1,
  cells = {font=\small},
  row{1} = {font=\small\bfseries},
  hborder{2-Y} = {pagebreak=no},
  rowsep = 3pt,
}
\toprule
Modelo & Uso & Licencia \\
\midrule
Granite 3.3 2B & Generación de respuestas en el perfil de menor consumo. & Apache 2.0~\cite{licencia_granite} \\
Gemma 4 E4B & Generación de respuestas y corrección de texto en el perfil de mayor capacidad. & Apache 2.0~\cite{licencia_gemma} \\
Multilingual E5 Small & Representación del texto para la búsqueda semántica. & \acs{MIT}~\cite{licencia_e5} \\
\bottomrule
\end{longtblr}

\noindent\textbf{Herramientas de desarrollo}\par

Las herramientas utilizadas para construir, probar o desplegar el proyecto
tampoco tienen por qué compartir una licencia. TypeScript utiliza Apache 2.0,
mientras que Vite se distribuye bajo \acs{MIT}. El código de Docker Compose utiliza
Apache 2.0, aunque Docker Desktop y los servicios alojados tienen sus propias
condiciones de uso~\cite{licencia_typescript,licencia_vite,licencia_compose,
docker_condiciones}. Terraform 1.15.8 emplea la
\textit{Business Source License} (\acs{BSL})\acused{BSL},
con un permiso adicional que limita determinados usos
competitivos, mientras que el proveedor de Cloudflare 5.23.0 se publica bajo
Apache 2.0~\cite{licencia_terraform,licencia_cloudflare_provider}. Estas
condiciones afectan a las herramientas, sin convertirse automáticamente en
la licencia del código que se desarrolla con ellas.

La variante de procesamiento con una \ac{GPU} de NVIDIA incorpora, además,
componentes sujetos a los acuerdos de \acs{CUDA} y \acs{cuDNN} (del inglés
\textit{CUDA Deep Neural Network})\acused{cuDNN}. Estos acuerdos permiten
utilizar las bibliotecas y redistribuir determinadas partes bajo sus
condiciones, por lo que deben revisarse los componentes incluidos cuando se
prepare una imagen que las contenga~\cite{nvidia_cuda_licencia,
nvidia_cudnn_licencia}.

\noindent\textbf{Servicios externos}\par

La búsqueda de juegos utiliza el catálogo de \textit{BoardGameGeek} y muestra
la atribución \guillemotleft{}Powered by BoardGameGeek\guillemotright{} cuando se consulta. Para aplicaciones
no comerciales, el proveedor permite normalmente el acceso gratuito,
sujeto al registro y la autorización de uso de su interfaz de programación de aplicaciones
(\acs{API}, del inglés \textit{Application Programming Interface})\acused{API}
y a la inclusión
del logo de atribución enlazado~\cite{boardgamegeek_api_condiciones}.

Para entregar los correos de verificación y recuperación de cuentas en el
despliegue público se utiliza Resend. A diferencia de Mailpit, que recoge
los mensajes en el entorno local, este proveedor recibe la dirección del
destinatario y el contenido del correo, que incluye el nombre de usuario
y el enlace temporal para completar la operación. No recibe los manuales
ni las conversaciones mediante este envío. Su uso se rige por las
condiciones del servicio, no por la licencia del código de
\textit{Manualito}~\cite{resend_condiciones}.

\noindent\textbf{Tratamiento de datos personales}\par

El envío de estos mensajes incorpora a Resend como encargado del
tratamiento. Su acuerdo de tratamiento de datos establece las condiciones
del servicio y las cláusulas contractuales tipo para las transferencias
internacionales~\cite{resend_tratamiento}. Según la información del proveedor,
los mensajes y registros se almacenan en Estados Unidos, aunque se elija
una región europea para enviarlos. Los planes Free, Pro y Scale conservan
los correos y registros durante 30 días, mientras que Enterprise permite
acordar otros plazos~\cite{resend_privacidad}.


Al compartir un manual, el usuario puede decidir si muestra su nombre o
aparece como \guillemotleft{}Anónimo\guillemotright{}, que es la opción por
defecto. Esta elección oculta el nombre en el manual y en las fuentes de
las respuestas, pero mantiene la cuenta asociada para comprobar la
propiedad y los permisos. Tampoco elimina datos que puedan estar escritos
en las páginas. Por ello, no constituye una anonimización completa de
los datos personales~\cite{aepd_datos_anonimizados} ni modifica los derechos
necesarios para subir y compartir el documento.

\subsubsection{Documentación}

La memoria y sus anexos se han redactado en \LaTeX{} y se distribuyen
bajo la licencia \acs{MIT}, al igual que el código y la documentación
del repositorio de \textit{Manualito}. Esta licencia permite consultar,
modificar y compartir la documentación siempre que se conserven el aviso
de autoría y el texto de la licencia~\cite{manualito_licencia,licencia_mit}.

\subsubsection{Imágenes y otros recursos}

Las tipografías de la interfaz, entre ellas Inter, Manrope y JetBrains Mono,
se distribuyen bajo la \textit{Open Font License} (\acs{OFL})\acused{OFL}.
Esta licencia permite incluirlas
en la aplicación conservando sus avisos y no se aplica al contenido escrito
con ellas. Los iconos de Lucide utilizan la licencia \acs{ISC} y conservan
también avisos \acs{MIT} de los iconos procedentes de Feather, de modo que sus
atribuciones deben mantenerse junto a los recursos que se distribuyan
\cite{licencia_ofl,licencia_lucide}.

Los recursos propios se publican bajo la licencia \acs{MIT}, pero los logotipos
y otros elementos de terceros conservan sus propias condiciones de uso,
aunque aparezcan en los diagramas o las capturas de la memoria.

Los usuarios deben disponer de los derechos necesarios sobre los manuales que
suban y compartan, sin perjuicio de las obligaciones legales de quien gestione
el servicio.

\subsubsection{Resumen}

La tabla~\ref{tabla:licencias-entrega} reúne las condiciones de los materiales
propios que acompañan a la entrega, a los que se añaden las licencias de
terceros descritas anteriormente.

\begin{longtblr}[
  caption = {Condiciones de distribución de los materiales del proyecto.},
  label = {tabla:licencias-entrega},
]{
  width = \linewidth,
  colspec = {X[3,l] X[3,l] X[4,l]},
  rowhead = 1,
  cells = {font=\small},
  row{1} = {font=\small\bfseries},
  hborder{2-Y} = {pagebreak=no},
  rowsep = 3pt,
}
\toprule
Material & Licencia o régimen & Condición principal \\
\midrule
Código propio & \acs{MIT} & Conservar el aviso de autoría y la licencia. \\
Documentación del programa & \acs{MIT} & Conservar el aviso de autoría y la licencia. \\
Memoria, diagramas y capturas propios & \acs{MIT} & Conservar el aviso de autoría y la licencia, respetando los elementos de terceros incluidos. \\
\bottomrule
\end{longtblr}

\endgroup
